\documentclass{rhelder3cv}

\author{Russell Wright Helder}
\title{PhD Candidate}
\affiliation{%
  Joint Program in Ancient Philosophy \\
  Department of Philosophy \\
  University of California, Berkeley%
}
\contact{%
  rwh@russellhelder.net \\
  \textdigits{(913) 428-0466}%
}
\date{August 12, \textdigits{2026}}

\begin{document}
\maketitle

\section{Education}

\begin{description}
  \item[PhD Candidate, Philosophy] \textacro{UC} Berkeley,
    \textdigits{2019}--present:
    to be completed by December \textdigits{2026}
  \item[MA with Distinction, Philosophy] Georgia State University,
    \textdigits{2019}
  \item[BSA, Chemistry \& BA, Classics] The University of Texas at Austin,
    \textdigits{2016}
\end{description}

\section{Areas of Research}

\begin{description}
  \item[Areas of Specialization] Ancient Greek philosophy,
    especially Aristotle’s ethics and moral psychology
  \item[Areas of Competence] Early modern philosophy,
    especially eighteenth century moral philosophy;
    contemporary ethics;
    philosophy of science
\end{description}

\section{Languages}

\begin{itemize}
  \item Ancient Greek (teaching experience)
  \item Latin (tutoring experience)
\end{itemize}

\section{Teaching}

\subsection{%
  \texorpdfstring{%
    Graduate Student Instructor
    \textnormal{(}\textacro{UC} Berkeley\textnormal{)}%
  }{%
    Graduate Student Instructor (UC Berkeley)%
  }%
}

\begin{enumerate}
  \item Philosophy \textdigits{100}: Philosophical Methods
    (Spring \textdigits{2026})

    \emph{%
      \small
      Intended to improve the student’s ability to read and write philosophy,
      with special emphasis placed on developing analytic skills.
      Students meet one-on-one with the Graduate Student Instructor every week
      to discuss the Instructor's feedback on their written assignments.%
    }
  \item Philosophy \textdigits{100}
    (Fall \textdigits{2025})
  \item Philosophy 2: Individual Morality and Social Justice
    (Spring \textdigits{2025})

    \emph{%
      \small
      An introduction to some central issues in moral and political philosophy.
      Students meet with the Graduate Student Instructor every week
      in small discussion sections.%
    }
  \item Philosophy \textdigits{25A}: Ancient Philosophy
    (Fall \textdigits{2024})

    \emph{%
      \small
      An introduction to ancient Greek philosophy,
      focusing mainly on the thought of Socrates, Plato, and Aristotle.
      Students meet with the Graduate Student Instructor every week
      in small discussion sections.%
    }
  \item Philosophy \textdigits{25B}: Modern Philosophy
    (Spring \textdigits{2023})

    \emph{%
      \small
      An introduction to 17th- and 18th-century philosophy,
      including Descartes, Elisabeth, Spinoza, Locke, Conway, Leibniz, Hume,
      and Kant.
      Students meet with the Graduate Student Instructor every week
      in small discussion sections.%
    }
  \item Philosophy \textdigits{100}
    (Fall \textdigits{2022})
  \item Greek 15: The Greek Workshop, Part I
    (Summer \textdigits{2022})

    \emph{%
      \small
      The first six weeks of a ten-week intensive introductory language course
      for students with little or no previous experience in Ancient Greek.
      In the first six weeks,
      students cover both semesters of elementary Greek.
      Students meet with the Graduate Student Instructor daily
      in small review sessions.%
    }
  \item Philosophy \textdigits{161}: Aristotle
    (Spring \textdigits{2022})

    \emph{%
      \small
      A concentrated study of Aristotle’s philosophy,
      including logic, theory of explanation, natural philosophy, psychology,
      metaphysics, and ethics.
      Students learn how to read Aristotle’s dense prose
      and form their own judgments
      about controversial matters of interpretation.
      Students meet with the Graduate Student Instructor every week
      in small discussion sections.%
    }
  \item Philosophy \textdigits{25A}
    (Fall \textdigits{2021})
  \item Philosophy \textdigits{25A}
    (Summer \textdigits{2021})
  \item Philosophy \textdigits{100}
    (Spring \textdigits{2021})
  \item Philosophy \textdigits{25A}
    (Fall \textdigits{2020})
\end{enumerate}

\subsection{%
  Instructor of Record
  \textnormal{(}Georgia State University\textnormal{)}%
}

\begin{enumerate}
  \item Philosophy \textdigits{1010}: Critical Thinking
    (Spring \textdigits{2019})

    \emph{%
      \small
      Aims to help students improve their critical thinking skills
      of identifying and evaluating arguments,
      and to thereby help students succeed in their other college courses.
      Instruction focuses on leading students through in-class exercises,
      preparing students for written \enquote{Standardize \& Evaluate} exams.%
    }
  \item Philosophy \textdigits{1010}
    (Fall \textdigits{2018})
  \item Philosophy \textdigits{1010}
    (Summer \textdigits{2018})
\end{enumerate}

\subsection{Other Teaching Positions}

\begin{enumerate}
  \item Tutor
    (Fall \textdigits{2016}--Spring \textdigits{2018})

    \emph{%
      \small
      Tutored middle-school students,
      high-school students,
      college students,
      and postgraduate students
      in Latin, chemistry, and academic writing,
      as well as general study skills
      and \textacro{ACT} preparation.%
    }
  \item Writing Across the Curriculum consultant
    for Philosophy \textdigits{4030}:
    Piety and the Gods in Ancient Philosophy
    (Georgia State University, Spring \textdigits{2018})

    \emph{%
      \small
      Gave feedback on rough drafts of all student papers.%
    }
  \item Supplemental Instruction leader for Philosophy \textdigits{1010}:
    Critical Thinking
    (Georgia State University, Fall \textdigits{2017})

    \emph{%
      Led small, student-led review sessions for Philosophy \textdigits{1010}.%
    }
\end{enumerate}

\section{Service}

\begin{description}
  \item[Equity Task Force] Department of Philosophy,
    \textacro{UC} Berkeley, \textdigits{2019}--\textdigits{2021}.
    Researched and co-authored
    \enquote{Best Practices Guide for Respectful Classroom Participation}
\end{description}

\section{Fellowships and Awards}

\begin{description}
  \item[Ralf F. Munster Fellowship] Department of Philosophy,
    Georgia State University, \textdigits{2018}.
    Awarded to an outstanding first-year graduate student
\end{description}

\section{Academic Presentations}

\subsection{%
  \texorpdfstring{%
    Plato, \texttitle{Meno} \textdigits{70a1}--\textdigits{73c5}%
  }{%
    Plato, Meno 70a1--73c5%
  }%
}

\begin{itemize}
  \item Invited.
    Seminar on Plato's \texttitle{Meno},
    Yale University
    (June 9--13, \textdigits{2025}).
    Co-presented with Verity Harte
\end{itemize}

\subsection{%
  \texorpdfstring{%
    Practical \textalt{Nous} of Universal Minor Premises:
    A New Interpretation of \texttitle{Nicomachean Ethics} \textdigits{6.11}%
  }{%
    Practical Nous of Universal Minor Premises:
    A New Interpretation of Nicomachean Ethics 6.11%
  }%
}

\begin{itemize}
  \item The Fifth Annual Rackham
    Interdisciplinary Workshop in Ancient Philosophy,
    University of Michigan, Ann Arbor
    (April 12, \textdigits{2024})
  \item American Philosophical Association,
    Pacific Division Meeting,
    San Francisco, CA
    (April 9, \textdigits{2023})
\end{itemize}

\subsection{%
  \texorpdfstring{%
    \texttitle{De Motu Animalium} 11: Involuntary and Non-Voluntary Movement%
  }{%
    De Motu Animalium 11: Involuntary and Non-Voluntary Movement%
  }%
}

\begin{itemize}
  \item Invited.
    Seminar on Aristotle's \textalt{De motu animalium},
    Yale University
    (June 12--16, \textdigits{2023}).
    Co-presented with \initials{M.M.} McCabe
\end{itemize}

\subsection{%
  A Bridge to Nowhere: Ensembles in Equilibrium Statistical Mechanics%
}

\begin{itemize}
  \item The 22nd Annual Conference
    of the International Society for the Philosophy of Chemistry,
    University of Bristol, UK
    (July 18, \textdigits{2018})
\end{itemize}

\subsection{Virtue and Wealth in Adam Smith: What’s the Difference?}

\begin{itemize}
  \item Midsouth Philosophy Conference,
    Rhodes College, Memphis, TN
    (March 23, \textdigits{2018})
\end{itemize}

\end{document}
